% Generated draft: requires mathematical and source review.
\documentclass[11pt]{article}
\usepackage[margin=25mm]{geometry}
\usepackage{fontspec}
\setmainfont{Latin Modern Roman}
\usepackage{amsmath,amssymb,mathtools,unicode-math}
\setmathfont{Latin Modern Math}
\usepackage{enumitem,needspace,xurl,hyperref}
\setlength{\parindent}{0pt}
\setlength{\parskip}{5pt}
\setlength{\emergencystretch}{3em}
\newcommand{\N}{\mathbb{N}}
\newcommand{\Z}{\mathbb{Z}}
\newcommand{\Q}{\mathbb{Q}}
\newcommand{\R}{\mathbb{R}}
\newcommand{\C}{\mathbb{C}}
\newcommand{\F}{\mathbb{F}}
\newcommand{\A}{\mathbb{A}}
\newcommand{\PP}{\mathbb{P}}
\newcommand{\E}{\mathbb{E}}
\newcommand{\eps}{\varepsilon}
\newcommand{\dd}{\,\mathrm{d}}
\begin{document}
\section*{UnsolvedMath ID: 20000449}
\textbf{Source problem code:} catalogue entry
\subsection{Definitions and mathematical statement}
Let \(k\) be \(\Q\) or a \(p\)-adic field \(\Q_p\), let \(n\geq2\), and let \(F\in k[x_1,\ldots,x_n]\) be a nonzero homogeneous cubic. Put \(X_F=V(F)\subset\PP^{n-1}_k\). A \(k\)-rational line on \(X_F\) is \(\PP(U)\) for a two-dimensional \(k\)-linear subspace \(U\subset k^n\) on which \(F\) vanishes identically. Equivalently, it is a \(k\)-point of the scheme of lines \(F_1(X_F)\subset\operatorname{Gr}(2,n)\). The original question does not impose smoothness, so it must not be imposed silently on the extremal problem.

Wooley's Question 16 asks for cubic hypersurfaces without a rational line in as many variables as possible:
\begin{enumerate}
\item over \(\Q_p\), with the dependence on the prime \(p\) made explicit;
\item over \(\Q\).
\end{enumerate}
One precise extremal formulation is to determine, for each such field,
\[
M_k=\sup\{n\geq2:\exists\,0\ne F\in k[x_1,\ldots,x_n]_3
\text{ with }F_1(X_F)(k)=\varnothing\}.
\]
Here \(k[x_1,\ldots,x_n]_3\) denotes degree-three homogeneous forms; finite upper bounds would make this supremum a maximum. The workshop also asks, in Colliot-Th\'el\`ene's formulation, for examples with \(X_F(k)\ne\varnothing\) but \(F_1(X_F)(k)=\varnothing\). This is an additional pointed requirement, not implied by absence of a line.

The opening request concerns local--global solubility for lines over \(\Q\): determine when
\[
F_1(X_F)(\Q_v)\ne\varnothing\ \text{for every place }v\text{ of }\Q
\quad\Longrightarrow\quad F_1(X_F)(\Q)\ne\varnothing.
\]
This implication is separate from the unrestricted maximum \(M_\Q\); a cubic without even a \(\Q_p\)-line is not a counterexample to that implication. The supplied aid reports that Jahnel and Loughran already produced failures for smooth cubic surfaces, while the extremal variable counts remain open.

The historical remarks say that point-free cubic forms exist in at most nine variables, so examples in at least ten variables would go beyond that source of line-free examples. They also assert a local fourteen-variable existence bound and tentatively write a global bound of ``37?''. The supplied verification explicitly treats the fourteen-variable statement as unverified and says the question mark after \(37\) is genuine; it reports that the global thirty-seven-variable assertion subsequently became a theorem of Wooley. These historical numbers are not adopted as verified present bounds here.

The attached aid reports, in its conventions, the more recent bounds
\[
11\leq M_\Q\leq30,\qquad
11\leq M_{\Q_p}\leq21\quad(p=2,3,5),
\]
and a smooth eleven-variable form over \(\Q\), called the Brandes--Dietmann form \(F_2\), having the rational point \(e_4=[0:0:0:1:0:\cdots:0]\) but no \(\Q_2\)-line and hence no \(\Q\)-line. Its coefficients are not supplied, so this record does not itself give an explicit construction. The aid further describes a three-layer discrete-valuation obstruction for an expression \(H+\pi G_1+\pi^2G_2\), using a line-free reduction of \(H\) and anisotropic reductions of \(G_1,G_2\). The variable blocks, complete hypotheses, and proof are absent; no unrestricted theorem about all such expressions is asserted in this draft.
\subsection{Short English statement}
Find the largest possible number of variables in a cubic form with no rational line over \(\Q\) or \(\Q_p\), including examples that still have a rational point. Separately, study the failure of the local--global principle for lines. Historical and attached research bounds are identified by their supplied level of verification.
\subsection{Sources}
\begin{itemize}
\item Wooley, Question 16, with the additional pointed question attributed to Colliot-Th\'el\`ene, AIM workshop \emph{Rational and integral points on higher dimensional varieties}: \url{https://aimath.org/WWN/qptsurface2/qptsurface2.pdf}.
\item Local registry \texttt{data/unsolvedmath/problems.json}, permanent ID \(20000449\), exact source code \texttt{catalogue entry}; canonical \texttt{aim-algebraic-number-theory-notes.json}, notes[100].
\item The supplied verification attributes a later thirty-seven-variable theorem to Wooley. The attached aid names Jahnel--Loughran for local--global counterexamples and Brandes--Dietmann for the form \(F_2\), without full bibliographic citations. Research attempt 1 is explicitly machine-generated, not peer reviewed, and requires independent verification before citation.
\end{itemize}
\end{document}
